@@M0@@
we have to show that the Modules @@M1@@ and @@M2@@
are canonically isomorphic. On the one hand, we have the canonical isomorphisms:
@@M3@@
On the other hand:
@@M4@@
Finally, it follows from EGA II 2.2.2, translated into sheaf language, that we have the canonical isomorphisms:
@@M5@@
whence, finally, the desired isomorphism.

It is clear that, if @@M6@@ is quasi-coherent, the Module @@M7@@ thus constructed is indeed the one defined in EGA II.

\sgaStableTarget{sga6.e06.para.2-3-3}2.3.3. It follows from 2.3.1 and 2.3.2 that the functor @@M8@@ is exact. Returning to the notation of 2.2, by passage to the derived categories we therefore deduce a functor still denoted @@M9@@:
@@M10@@

We also see from the preceding constructions that the functor
